\chapter{Launching}\label{ch:launching}

\section{Staffing}\label{sec:staffing}

Nella fase di planning si è stimato essere necessario avere 5 sviluppatori full-time per completare il progetto. Per questo motivo, il team che ha preso parte alla fase di planning corrisponde con il developer team: Asia, Giacomo, Linda (che ha anche ricoperto il ruolo di product owner), Nicolas e Nicolò.
Tutti i membri del team hanno già grande esperienza con progetti Scrum, per questo motivo non è stato necessario ingaggiare uno Scrum Master che guidi il team nell'applicazione di questa metodologia.

I membri del team sono ``persone fatte a T'': con questo termine si indica una persona con capacità di comprendere e affrontare uno spettro ampio di temi e problematiche -- rappresentato dalla barra orizzontale della T -- e allo stesso tempo con una forte competenza in una specifica materia -- rappresentata dalla barra verticale. Questo permette al team di avere una grande varietà di competenze e di poter affrontare qualsiasi problema che si presenti durante lo sviluppo del progetto. In seguito sono riportati i punti di forza di ciascuno sviluppatore:
\begin{itemize}
  \item Asia (\emph{accommodating}): ha grande esperienza nello sviluppo frontend, di interfacce utente con particolare attenzione alla user experience
  \item Giacomo (\emph{assimilating}): è esperto nella realizzazione di sistemi basati su una architettura funzionale e in linguaggi di programmazione funzionali 
  \item Linda (\emph{accommodating}): è esperta di Domain Driven Design e di modellazione di domini complessi
  \item Nicolas (\emph{diverging}): ha grande esperienza nell'utilizzare pratiche agili e strumenti automatizzati per garantire la qualità del codice, la scalabilità del sistema e la sicurezza delle operazioni
  \item Nicolò (\emph{converging}): è esperto nella progettazione, training e valutazione di sistemi di machine learning
\end{itemize}

\begin{figure}
  \centering
  \begin{tikzpicture}
    \draw[line width=30pt] (0,0) -- (0,4) node [midway, rotate=90, color=white] {conoscenze};  % linea verticale
    \draw[line width=30pt] (-2,4) -- (2,4) node[midway, color=white] {competenze}; % linea orizzontale
  \end{tikzpicture}
  \caption{Persona a T}
  \label{fig:croce}
\end{figure}

\todo[inline]{Decidere se fare o meno il Project Definition Statement}

\section{Definizione delle regole operative}
\subsection{Project meeting}
È stata definita la gamma di possibili meeting a cui il team dovrà partecipare, la loro durata e cadenza.

\subsubsection{Standup meeting}
Gli \emph{standup meeting} si svolgono con cadenza giornaliera e non impiegano più di 10 minuti. Questi meeting coinvolgono tutti i membri in piedi ogni mattina prima di iniziare la giornata.
Questi si svolgono in maniera informale e non hanno una scaletta esplicita: i membri del team di sviluppo discutono su ciò che stanno facendo e presentano gli eventuali problemi riscontrati. 
Nel caso in cui siano presenti problemi che impediscano di proseguire con lo sviluppo, il team si riunirà in un \emph{problem resolution meeting} per discutere più nel dettaglio del problema.

\subsubsection{Problem Resolution Meeting}
I problem resolution meeting vengono indetti in maniera straordinaria qualora durante uno standup meeting emerga un problema grave. Alla riunione partecipano il product owner, il \emph{problem owner} ed eventuali sviluppatori che possono essere affetti dal problema. Le riunioni si svolgono secondo una scaletta predefinita:
\begin{itemize}
  \item Vengono raccolte informazioni sul problema
  \item Si cercano possibili soluzioni tramite \emph{brainstorming}
  \item I partecipanti si consultano per decidere la soluzione migliore; nel caso in cui sia necessario potrebbero essere interpellati anche altri membri del team per chiedere un loro parere
  \item Si aggiorna lo \emph{sprint backlog} in modo da includere i task necessari a implementare la soluzione scelta al problema
\end{itemize}

\subsubsection{Sprint review}

\subsubsection{Sprint retrospective}

\subsubsection{Sprint planning}


